%%%%%%%%%%%%%%%%%
% This is an sample CV template created using altacv.cls
% (v1.1.5, 1 December 2018) written by LianTze Lim (liantze@gmail.com). Now compiles with pdfLaTeX, XeLaTeX and LuaLaTeX.
%
%% It may be distributed and/or modified under the
%% conditions of the LaTeX Project Public License, either version 1.3
%% of this license or (at your option) any later version.
%% The latest version of this license is in
%%    http://www.latex-project.org/lppl.txt
%% and version 1.3 or later is part of all distributions of LaTeX
%% version 2003/12/01 or later.
%%%%%%%%%%%%%%%%

%% If you need to pass whatever options to xcolor
\PassOptionsToPackage{dvipsnames}{xcolor}

%% If you are using \orcid or academicons
%% icons, make sure you have the academicons
%% option here, and compile with XeLaTeX
%% or LuaLaTeX.
% \documentclass[10pt,a4paper,academicons]{altacv}

%% Use the "normalphoto" option if you want a normal photo instead of cropped to a circle
% \documentclass[10pt,a4paper,normalphoto]{altacv}

\documentclass[10pt,letter,ragged2e]{altacv}
\usepackage{hyperref}
%% AltaCV uses the fontawesome and academicon fonts
%% and packages.
%% See texdoc.net/pkg/fontawecome and http://texdoc.net/pkg/academicons for full list of symbols. You MUST compile with XeLaTeX or LuaLaTeX if you want to use academicons.

% Change the page layout if you need to
\geometry{left=1cm,right=9cm,marginparwidth=6.8cm,marginparsep=1.2cm,top=1.25cm,bottom=1.25cm}

% Change the font if you want to, depending on whether
% you're using pdflatex or xelatex/lualatex
\ifxetexorluatex
  % If using xelatex or lualatex:
  \setmainfont{Lato}
\else
  % If using pdflatex:
  \usepackage[utf8]{inputenc}
  \usepackage[T1]{fontenc}
  \usepackage[default]{lato}
\fi

% Change the colours if you want to
\definecolor{Orange}{HTML}{AC3C00}
\definecolor{DarkOrange}{HTML}{802d00}
\definecolor{SlateGrey}{HTML}{2E2E2E}
\definecolor{LightGrey}{HTML}{666666}
\colorlet{heading}{DarkOrange}
\colorlet{accent}{DarkOrange}
\colorlet{emphasis}{SlateGrey}
\colorlet{body}{LightGrey}

% Change the bullets for itemize and rating marker
% for \cvskill if you want to
\renewcommand{\itemmarker}{{\small\textbullet}}
\renewcommand{\ratingmarker}{\faCircle}

%% sample.bib contains your publications
\addbibresource{sample.bib}

\begin{document}
\name{Charles Rumford}
\tagline{Security and automation focused network architect with a love of usablity}
\personalinfo{%
  % Not all of these are required!
  % You can add your own with \printinfo{symbol}{detail}
    \email{\href{mailto:charlesr@tallwireless.com}{charlesr@tallwireless.com}}
  \phone{+1 267.702.6721/+44 77 628 22768}
  %% You MUST add the academicons option to \documentclass, then compile with LuaLaTeX or XeLaTeX, if you want to use \orcid or other academicons commands.
  % \orcid{orcid.org/0000-0000-0000-0000}
}

%% Make the header extend all the way to the right, if you want.
\begin{fullwidth}
\makecvheader
\end{fullwidth}

%% Depending on your tastes, you may want to make fonts of itemize environments slightly smaller
% \AtBeginEnvironment{itemize}{\small}

%% Provide the file name containing the sidebar contents as an optional parameter to \cvsection.
%% You can always just use \marginpar{...} if you do
%% not need to align the top of the contents to any
%% \cvsection title in the "main" bar.
\cvsection[resume-p1sidebar]{Experience}

\cvevent{Principal Infrastructure Architect}{Summit Hosting}{Jan 2021 - Ongoing}{}
\begin{itemize}
\item Development and deployment of fully automated EVPN fabrics to support customer traffic
\item Lead efforts to advance network automation effort organization wide
\item Research, development, and deployment of new networking services including automated DDoS
    scrubbing and multi-tenent load balancing services
\item Serve as final escalation point for all networking related services
\end{itemize}

\divider

\cvevent{Network Architect}{University of Pennsylvania}{Oct 2012 -- Feb 2021}{}
\begin{itemize}
    \item Former positions: Network Engineer (2012-2016), and Senior
        Network Engineer (2016-2019)
\item Design and deployment of first university border firewall and IPS
    resulting in blocking millions of potential attacks per month
\item Strategy development and deployment of large scale network architecture
    changes to support network segmentation efforts with MPLS
\item Deployment of carrier grade NAT to free 64,000 public IP addresses
    saving the university \$1 million in additional IP address costs
\item Advocate for the use of IPv6 addressing for all services on the network
\item Serve as final escalation point for all networking related services
\end{itemize}

%\divider
%
%\cvevent{Senior Network Engineer}{University of Pennslyvania}{Jan 2016 -- Jan 2019}{Philadeplia}
%\begin{itemize}
%\item goobar
%\end{itemize}

\divider

\cvevent{Adjunct Professor}{Drexel University}{Jan 2014 - Sept 2016}{}
\begin{itemize}
\item Taught courses in C programming and Linux system administration
\item Development of lectures, course work, and assessments
\end{itemize}

\divider


\cvevent{Network Analyst}{Drexel University}{Aug 2010 -- Sept 2012}{}
\begin{itemize}
\item Primarily responsible for the management of the wireless infrastructure
    supporting 20,000 connected devices
\item Deployment of new wireless authentication systems resulting in improved
    reliability and visibility into authentications
\item development and deployment of AP deployment scheme for residence halls
    to improve coverage for students
\end{itemize}

%  \divider
%
%  \cvevent{Technical Director}{Drexel University Computing Academy}{June 2006 -- July 2014}{}
%
%  \divider
%
%  \cvevent{Linux Systems Administrator}{Drexel University}{Sept 2006 - Dec 2008}{}
%  \begin{itemize}
%  \item Strategy development and deployment of system wide monitoring and notification system
%  \item Advocate for the migration from Gentoo to Ubuntu as the primary
%      distribution to reduce management overhead
%  \item Maintenance of centralize authentication and file services using LDAP and
%      NFS
%  \end{itemize}

%\cvsection{Projects}
%
%\cvevent{Project 1}{Funding agency/institution}{}{}
%\begin{itemize}
%\item Details
%\end{itemize}
%
%\divider
%
%\cvevent{Project 2}{Funding agency/institution}{Project duration}{}
%A short abstract would also work.
%
%\medskip
%
%\cvsection{A Day of My Life}
%
%% Adapted from @Jake's answer from http://tex.stackexchange.com/a/82729/226
%% \wheelchart{outer radius}{inner radius}{
%% comma-separated list of value/text width/color/detail}
%\wheelchart{1.5cm}{0.5cm}{%
%  6/8em/accent!30/{Sleep,\\beautiful sleep},
%  3/8em/accent!40/Hopeful novelist by night,
%  8/8em/accent!60/Daytime job,
%  2/10em/accent/Sports and relaxation,
%  5/6em/accent!20/Spending time with family
%}
%
%\clearpage
%\cvsection[page2sidebar]{Publications}
%
%\nocite{*}
%
%\printbibliography[heading=pubtype,title={\printinfo{\faBook}{Books}},type=book]
%
%\divider
%
%\printbibliography[heading=pubtype,title={\printinfo{\faFileTextO}{Journal Articles}},type=article]
%
%\divider
%
%\printbibliography[heading=pubtype,title={\printinfo{\faGroup}{Conference Proceedings}},type=inproceedings]

%% If the NEXT page doesn't start with a \cvsection but you'd
%% still like to add a sidebar, then use this command on THIS
%% page to add it. The optional argument lets you pull up the
%% sidebar a bit so that it looks aligned with the top of the
%% main column.
% \addnextpagesidebar[-1ex]{page3sidebar}


\end{document}
